\begin{theorem}[Bernoulli’s inequality]
Let \(x\in\mathbb{R}\) satisfy \(x\ge -1\) and let \(n\in\mathbb{Z}\) with \(n\ge 0\). Then  
\[
(1+x)^{n}\;\ge\;1+nx .
\]
\end{theorem}

\begin{proof}
We proceed by induction on the integer exponent \(n\).

\medskip\noindent\textbf{1. Base cases.}
\begin{itemize}
\item \textbf{\(n=0\).}  By convention we define \(0^{0}=1\); this agrees with the limit
\(\displaystyle \lim_{y\to0^{+}}y^{0}=1\). Hence
\begin{align*}
(1+x)^{0}=1,\qquad 1+0\cdot x = 1 .
\end{align*}
The inequality holds with equality for every admissible \(x\ge -1\).

\item \textbf{\(n=1\).}  Direct computation gives
\begin{align*}
(1+x)^{1}=1+x,\qquad 1+1\cdot x = 1+x .
\end{align*}
Equality again holds for all \(x\ge -1\).

\item \textbf{\(x=-1\).}  Here \(1+x=0\). For any \(n\ge 0\)
\[
(1+x)^{n}=0^{n}=%
\begin{cases}
1, & n=0,\\[2pt]
0, & n\ge 1 .
\end{cases}
\]
The right–hand side equals \(1+n(-1)=1-n\). Thus
\[
0^{n}\ge 1-n .
\]
This is true because for \(n=0\) both sides are \(1\), and for \(n\ge 1\) the left side is \(0\) while the right side is at most \(0\). Equality occurs only when \(n=0\) or \(n=1\).
\end{itemize}
Hence the statement is valid for the initial exponents and for the special value \(x=-1\).

\medskip\noindent\textbf{2. Induction hypothesis.}
Assume that for a fixed integer \(k\ge 1\) the inequality
\begin{align}
(1+x)^{k}\;\ge\;1+kx \tag{IH}
\end{align}
holds for every real \(x\ge -1\).

\medskip\noindent\textbf{3. Multiplication by the non‑negative factor \(1+x\).}
Since \(x\ge -1\) we have \(1+x\ge 0\). Multiplying both sides of (IH) by this
non‑negative quantity preserves the direction of the inequality:
\begin{align}
(1+x)^{k+1}= (1+x)^{k}(1+x)\;\ge\;(1+kx)(1+x). \tag{1}
\end{align}

\medskip\noindent\textbf{4. Expansion of the right‑hand side.}
\begin{align}
(1+kx)(1+x)=1+(k+1)x+kx^{2}. \tag{2}
\end{align}

\medskip\noindent\textbf{5. Dropping a non‑negative term.}
For any real \(x\) we have \(x^{2}\ge 0\), and because \(k\ge 1\) the product
\(k x^{2}\) is non‑negative:
\begin{align}
kx^{2}\ge 0. \tag{3}
\end{align}
Consequently, from (2) we obtain
\begin{align}
(1+kx)(1+x)\;\ge\;1+(k+1)x . \tag{4}
\end{align}

\medskip\noindent\textbf{6. Inductive step.}
Combining (1) and (4) yields
\begin{align}
(1+x)^{k+1}\;\ge\;1+(k+1)x ,
\end{align}
which is precisely the desired inequality for the exponent \(k+1\).

\medskip\noindent\textbf{7. Completion of the induction.}
The property holds for \(n=0\) and \(n=1\); the inductive step shows that
truth for \(n=k\) implies truth for \(n=k+1\).  By the Principle of
Mathematical Induction, the inequality
\[
(1+x)^{n}\;\ge\;1+nx
\]
holds for every integer \(n\ge 0\) and every real \(x\ge -1\).

\medskip\noindent\textbf{8. Equality cases.}
Equality in the inductive step can arise in two ways:
\begin{enumerate}
\item The discarded non‑negative term in (5) vanishes, i.e. \(k x^{2}=0\).
Since \(k\ge 1\), this forces \(x=0\).
\item The factor \(1+x\) that multiplies both sides in (3) is zero, i.e.
\(x=-1\).  In this situation (1) is an equality irrespective of the
induction hypothesis.
\end{enumerate}
Collecting all possibilities, equality holds exactly in the following
situations:
\begin{itemize}
\item \(n=0\) – both sides equal \(1\) for any \(x\ge -1\);
\item \(n=1\) – both sides equal \(1+x\) for any \(x\ge -1\);
\item \(x=0\) – \((1+0)^{n}=1=1+n\cdot0\) for any \(n\ge 0\);
\item \(x=-1\) and \(n=1\) – both sides are \(0\).
\end{itemize}
For any \(n\ge 2\) and any \(x\neq 0,-1\) the term \(k x^{2}\) (and, more
generally, the higher‑order terms of the binomial expansion) is strictly
positive, making the left‑hand side strictly larger than \(1+nx\).  Hence no
other equality cases exist. \(\square\)
\end{proof}